% ----------------------------------------------------------
\chapter{Discussion}
\label{chap:discussion}
% ----------------------------------------------------------

\setlength{\epigraphwidth}{0.6\textwidth}
\epigraph{%
    \vspace*{-1cm}
    \itshape% 
    It's tough to make predictions, especially about the future.
}{Yogi Berra}

\section{Summary of key findings}

Increase in gating intensity tended to increase 
observation coverage, mean squared error, and data mismatch,
albeit with case-dependent sensitivities.
%
Localization yielded mixed results: for the Egg and real field cases,
effects on coverage and MSE were minimal, especially for intense enough gatings,
while the effect in the Olympus case was drastic.
%
Cumulative distribution plots of the field oil production total suggest 
that ground-truths are unbiased in the Egg and real field cases,
but somewhat biased for the Olympus case.
%
Mild gating intensities (2 in the Egg and Olympus cases, 4 in the real field case) produced similar results 
to the standard DSI-ESMDA for same localization configuration,
while small gating (1 for the Egg case, 1 and 2 for the Olympus and real field cases)
had worse coverages than the standard DSI-ESMDA.
 
\section{Interpretation of findings}

\subsection{OC, MSE and mismatch increase}
That the metrics employed tended to increase with the increase in gating intensity is not surprising,
as the objective of Gated DSI-ESMDA lies exactly in increasing posterior variance,
i.e., its spread.
%
This naturally leads to higher coverages at the cost of higher mismatch and MSE. 

As for the differing standard deviations on the bar plots of the Egg and Olympus cases,
they are most probably linked to ground-truth representativeness and reservoir complexity.
As indexed by FOPT range and distribution, 
the Egg case is very simple and the picked ground-truths, relatively well-represented.
On the other hand, the Olympus reservoir is much more complex, 
and picked groundtruths tended to be biased towards low or high oil productions relative to the prior ensemble.
Indeed, ground-truths such as 40 and 35 showed FOPT that was close only to a handful realizations.
As such, metrics on the Olympus case varied more, and on the Egg case, less. 

\subsection{Mixed effects of localization}
Neither the Egg case nor the real field case demonstrated much sensitivity towards localization of the Kalman-like gain,
but the Olympus case did.
The reason behind this result is interpreted as combination of five factors:
reservoir complexity (as indexed by FOPT), 
ground-truth representativeness, 
ensemble size,
data dimentionality,
feature selection, and 
localization tuning.

For the Egg case, this could be explained by the rationale that 
the reservoir is simple enough that the number of realizations contitute a rich enough sampling of the prior distribution.
Importantly, this is also the case study with the smallest amount of features 
(\cref{tab:simulation-parameters-features}),
and reduced dimentionality help with the estimation of covariances.
%
Lastly, given the narrow range of this ensemble's FOPT and the roughly uniform sampling of ground-truths,
it is fair to say no one ground-truth was left unrepresented, which explains the small standard deviations.
%
In short, the Egg reservoir is too well-behaved for localization to be effective.

Conversely, the Olympus reservoir is much more complex while also having a much smaller ensemble
and a lot more features, aggravating covariance estimation.
Localization was conceived exactly for such situations, so it is not suprising it was effective here.

As for the real case, reservoir complexity and data dimentionality peak,
but so does ensemble size.
%
And unlike the Olympus case, the real field should be well-represented, at least FOPT-wise.
Rather, the main reason for localization not being effective 
are more likely to be the adjustment of its parameters 
(anisotropy, critical lengths and time)
and feature selection.
%
Spacial localization, which was not performed, could have helped improve localization.
%
Notedly, feature selection by have been inadequate for this reservoir,
as indicated by the high MSE and mismatch values.
Cases like well I3 BHP in \cref{fig:figures-real_field}, 
where many if not all reference points are at the outskirts of the prior 
occurred multiple times,
massively impairing the search for an adequate posterior.

In that situation, methods that sequentially assimilate data points, 
such as the ensemble Kalman filter, might have an easier time matching data.
%
Still, the most adequate thing to do seems to be reselecting the data features used in prediction. 

\subsection{Mild and small gating intensities versus standard DSI-ESMDA}
It was commented (\cref{sec:gated-dsi-esmda}) 
that gating intensity reverts the posterior to the prior.
If the posterior given by the standard DSI-ESMDA is somewhere in that continuum,
then it stands to reason that Gated DSI-ESMDA could approximate it simply by tuning the gating intensity.

Conversely, small gatings might actually produce worse posterior spreads,
and exposed in the findings.
%
Since Gated DSI-ESMDA does away with observation data resampling,
small gatings possibly underestimate uncertainty (i.e., likelihood covariance)
at least when compared to standard DSI-ESMDA,
producing posteriors that are not well-matched to data. 
% In practice, knowing the values of such intensities gives a baseline on the performance of DSI-ESMDA itself,
% and could be used when deciding whether posterior variance should increase or decrease. 

\section{Evaluation of existing theories and models}

Findings are generally in agreement with the theory stablished in 
\cref{chap:theory-framework}, as discussed above.
%
While Gated DSI-ESMDA provides a straightforward, effective mechanism for 
increasing posterior covariance, too small a gating intensity could produce worse spreads than the standard DSI-ESMDA
as doing away with observation data resampling could underestimate uncertainty.
%
This is remediated as the intensity increases, and mild intensities were shown to produce 
similar results to DSI-ESMDA.
%
How large the intensity should be depends on how tolerable are the 
increase in MSE and data-mismatch that come as a trade-off.
